The quasilinear diffusion coefficients (\ref{dbnb}),
(\ref{drf0D_dke}) and (\ref{drf0F_dke}) describe the interaction
between electrons and a discrete set ($b$) of RF plane waves
(\ref{Ecomp}) with frequencies $\omega_{b}$. In addition, we can
linearly superpose a discrete set of plane waves with the same
frequency but different parallel wave number $k_{\Vert}$ because
of the linearity of the QL operator (\ref{qloporig}) in
$\mathbf{k}$. This set of plane waves is typically represented by
rays. Each ray is characterized by: the wave frequency
$\omega_{b}$, the poloidal location $\theta_{b}$, the index of
refraction $\mathbf{N}_{b}$, the polarization vector
$\mathbf{e}_{b}$ and the power flow $\mathbf{\Phi}_{b}$. In
general, the ray propagation path $\left(  r_{b},\theta_{b}\right)
$ and the evolution of the index of refraction $\mathbf{N}_{b}$
are determined by ray-tracing (RT) calculations. When the wave
properties are determined by RT calculations, the launched power
spectrum in $k_{\Vert}$ is typically decomposed in an array of
rays with different poloidal launching angle and/or $k_{\Vert}$
power spectrum. This distribution depends on the antenna and other
launching parameters. The evolution of each ray is determined
separately, by the RT code, and the contributions of each ray to
wave-particle interaction are separately accounted for in the sum
over $b$ in (\ref{dtensor1}-\ref{dtensor4}). The separation of
these contributions is justified if their $k_{\Vert}$ power
spectrums do not overlap.

We do not address the RT problem in this work. Therefore, we assume that the
path $\theta_{b}\left(  \psi\right)  $ and the parallel index of refraction
$N_{b\Vert}$ are determined either by coupling of DKE to a RT code, or by any
other propagation model. RT calculations require to calculate the dispersion
tensor and solve the dispersion relation, and therefore they can provide input
for the remaining wave properties $\left(  N_{b\bot},\mathbf{e}_{b}%
,\mathbf{\Phi}_{b}\right)  $. However, these properties can also be determined
by solving locally the wave equation, once $\psi,\theta_{b}\ $and $N_{\Vert}$
are known. Keeping the calculation of $\left(  N_{b\bot},\mathbf{e}%
_{b},\mathbf{\Phi}_{b}\right)  $ within DKE allows us to use a more advanced
model - such as the fully relativistic dispersion solver R2D2 - for the
evaluation of these important wave properties. A numerical code is required to
solve the wave equation, and evaluate the diffusion coefficient\ (\ref{dbnb})
in the general case. However, it is possible to obtain analytic expressions in
some cases, assuming that the waves can be treated in the cold plasma limit.
The calculation of $\overline{D}_{b,n}^{\text{RF}(0)}(p,\xi_{0})$ in the cold
plasma model is performed in Appendix \ref{cold}. In addition, simplified
expression and the comparison with operators used in the litterature are given
for lower-hybrid (LH) and electron-cyclotron (EC) waves in appendices
\ref{LHCD} and \ref{ECCD} respectively.

Note that the cold plasma model, if used in the RT calculations, does not
account for the wave-particle interaction and the power absorption from the
rays. The power absorbed from each ray and deposited on electrons can be
evaluated if a hot plasma model is used, but even then, the calculation is not
consistent with DKE because quasilinear effects are not accounted for. In
order to have a consistency between RT and DKE calculations, the power
absorption calculated by DKE must be inserted back in the RT calculations. In
the case where rays turn back toward the edge (as it can happen in LHCD),
iterations between RT and DKE calculations are necessary to ensure self-consistency.

\subsubsection{Modeling of RF $k_{\Vert}$ spectrum}

We consider two different models for the $k_{\Vert}$ spectrum of a given ray.

\paragraph{Square power spectrum in $k_{\Vert}$}

\bigskip Here, the power spectrum of a given ray is taken as being constant in
$N_{\Vert}$ between two limits $N_{b\Vert\min}$ and $N_{b\Vert\max}$.
Therefore, we operate the following transformation
\begin{equation}
\delta\left(  N_{\Vert}-N_{\Vert\text{res},b}\right)  \rightarrow\frac
{1}{\Delta N_{b\Vert}}H\left(  N_{\Vert\text{res},b}-N_{b\Vert\min}\right)
H\left(  N_{b\Vert\max}-N_{b\Vert\text{res},b}\right)  \label{square}%
\end{equation}
where%
\begin{equation}
\Delta N_{b\Vert}=N_{b\Vert\max}-N_{b\Vert\min}%
\end{equation}

In this case, the RT calculations and wave equation are solved for the central
value.%
\begin{equation}
N_{b\Vert,0}=\frac{N_{b\Vert\max}+N_{b\Vert\min}}{2}%
\end{equation}
which is a good approximation only when
\begin{equation}
\frac{\Delta_{b}N_{\Vert}}{N_{b\Vert,0}}\ll1
\end{equation}

For simplification and benchmarking purposes, in LHCD, one can consider only
one ray with a square spectrum. Then, the limits $N_{b\Vert\min}$ and
$N_{b\Vert\max}$ are respectively given by the accessibility condition and the
condition for strong linear Laudau damping.

\paragraph{Gaussian power spectrum in $k_{\Vert}$}

\bigskip The power spectrum of a given ray can be assumed to have a Gaussian
dependence on $N_{\Vert}$, centered around a value $N_{b\Vert,0}$ and with a
Gaussian width $\Delta N_{b\Vert}$. Then, we operate the transformation%
\begin{equation}
\delta\left(  N_{b\Vert}-N_{\Vert\text{res}}\right)  \rightarrow\frac{1}%
{\sqrt{\pi}\Delta N_{b\Vert}}\exp\left[  \frac{-\left(  N_{\Vert\text{res}%
}-N_{b\Vert,0}\right)  ^{2}}{\Delta N_{b\Vert}^{2}}\right]  \label{gaussian}%
\end{equation}

In this case too, the RT calculations and wave equation are solved for the
central value $N_{b\Vert,0}$, which is a good approximation only when
\begin{equation}
\frac{\Delta_{b}N_{\Vert}}{N_{b\Vert,0}}\ll1 \label{narrow}%
\end{equation}

For simplification and benchmarking purposes, in ECCD, one can consider only
one ray with a Gaussian spectrum. Then, the wave properties $N_{b\Vert,0}$ and
$\Delta N_{b\Vert}$ are determined by the beam characteristics, as described
in \ref{beam}.

\subsubsection{Beam size and power spectrum\label{beam}}

\bigskip If we consider a Gaussian beam of waist $d$, of negligible
diffraction angle, and of negligible dispersion, such that the wave vector is
in the direction of the beam, and $N\simeq1$, we get that by the properties of
Fourier transform, the width $\Delta N_{\Vert}$ is related to the beam waist
$d$ as
\begin{equation}
\Delta N_{\Vert}=\frac{N_{\bot}}{N}\frac{c}{\omega d}%
\end{equation}
This conditions are approximately satisfied for EC\ beamsif the density is low
($\omega_{pe}\ll\omega$).

Then, and we can verify that for a beam waist of a few centimeters ($d=5$ cm,
$f=110$ GHz, $N_{\Vert}\simeq0.2\rightarrow\Delta N_{\Vert}\simeq0.01$), the
condition (\ref{narrow}) is satisfied, and there is no nead to decompose the
beam spectrum into many rays.
